% GENERATED FILE. Edit canonical lessons, then run npm run generate:books.
% canonical-source-hash: fnv1a64:8d70ce51f77e53b6
% canonical-lessons: JA-C406-torii, JA-C406-kuguru, JA-C406-keidai, JA-C406-kannushi, JA-C406-sasageru

\chapter{At the Shrine Gate}
\label{ch:ja-at-the-shrine-gate}
\hlchaptermodality{406}

\begin{chapteropening}
\textbf{By the end of this chapter:} \emph{I can visit a shrine and follow its customs.}

Complete the last lesson of chapter 406: I can visit a shrine and follow its customs.
\end{chapteropening}

\section[\texorpdfstring{\ja{とりい} (torii)}{とりい (torii)}]{\ja{とりい} (torii) — a shrine gate}

\label{lesson:JA-C406-torii}



\begin{quote}
Before the new one: say the Japanese for to become rough, then the Japanese for to cool off.
\end{quote}

\subsection*{You'll want to know: \ja{とりい}}
\textbf{\ja{とりい}} — \emph{torii} — \textquotedblleft{}a shrine gate\textquotedblright{}.

\begin{quote}
\textbf{\ja{じんじゃの} \ja{まえに} \ja{あかい} \ja{とりいが} \ja{あります。}} — \emph{jinja no mae ni akai torii ga arimasu} — \textquotedblleft{}There is a red gate in front of the shrine.\textquotedblright{}
\end{quote}

\begin{cousinweb}
\textbf{\ja{とりい}} (\emph{torii}) is a native Japanese word (wago), inherited rather than borrowed from Chinese, spelled with the kanji for \textquotedblleft{}bird\textquotedblright{} and \textquotedblleft{}reside\textquotedblright{}; how the gate came by that name is not settled (JMdict, KANJIDIC2).
\end{cousinweb}

\begin{grammarlens}[title={What you've built}]
One more word for getting around.
\end{grammarlens}

\subsection*{Your turn}
\begin{itemize}
\raggedright
  \item \emph{Say it:} \emph{torii}
  \item \emph{Say it:} \emph{torii}, once more
  \item \emph{Say it:} \emph{jinja no mae ni akai torii ga arimasu} — There is a red gate in front of the shrine.
  \item \emph{Recall:} say the Japanese for to become rough, then the Japanese for to cool off, then say \emph{torii} again
\end{itemize}

\begin{tcolorbox}[breakable,colback=teal!4,colframe=teal!35!black,title={Before you move on}]
What does \ja{とりい} mean? (\textquotedblleft{}A shrine gate\textquotedblright{}.) Say it once more.
\end{tcolorbox}
\section[\texorpdfstring{\ja{くぐる} (kuguru)}{くぐる (kuguru)}]{\ja{くぐる} (kuguru) — to pass under; to go through}

\label{lesson:JA-C406-kuguru}



\begin{quote}
Before the new one: say the Japanese for to cool off, then the Japanese for a shrine gate.
\end{quote}

\subsection*{You'll want to know: \ja{くぐる}}
\textbf{\ja{くぐる}} — \emph{kuguru} — \textquotedblleft{}to pass under; to go through\textquotedblright{}.

\begin{quote}
\textbf{\ja{とりいを} \ja{くぐって、じんじゃに} \ja{はいりました。}} — \emph{torii o kugutte, jinja ni hairimashita} — \textquotedblleft{}We passed under the gate and went into the shrine.\textquotedblright{}
\end{quote}

Polite \textbf{\ja{くぐります}}, \emph{kugurimasu}; a godan verb. Past (polite): \textbf{\ja{かんこうきゃくが} \ja{とりいを} \ja{くぐりました。}}, \emph{kankōkyaku ga torii o kugurimashita}.

\begin{cousinweb}
\textbf{\ja{くぐる}} (\emph{kuguru}) is a native Japanese word (wago), inherited rather than borrowed from Chinese; its kanji was attached later as a kun'yomi (JMdict).
\end{cousinweb}

\begin{grammarlens}[title={What you've built}]
One more word for getting around.
\end{grammarlens}

\subsection*{Your turn}
\begin{itemize}
\raggedright
  \item \emph{Say it:} \emph{kuguru}
  \item \emph{Say it:} \emph{kuguru}, once more
  \item \emph{Say it:} \emph{torii o kugutte, jinja ni hairimashita} — We passed under the gate and went into the shrine.
  \item \emph{Recall:} say the Japanese for to cool off, then the Japanese for a shrine gate, then say \emph{kuguru} again
\end{itemize}

\begin{tcolorbox}[breakable,colback=teal!4,colframe=teal!35!black,title={Before you move on}]
What does \ja{くぐる} mean? (\textquotedblleft{}To pass under; to go through\textquotedblright{}.) Say it once more.
\end{tcolorbox}
\section[\texorpdfstring{\ja{けいだい} (keidai)}{けいだい (keidai)}]{\ja{けいだい} (keidai) — shrine grounds; temple grounds}

\label{lesson:JA-C406-keidai}



\begin{quote}
Before the new one: say the Japanese for a shrine gate, then the Japanese for to pass under.
\end{quote}

\subsection*{You'll want to know: \ja{けいだい}}
\textbf{\ja{けいだい}} — \emph{keidai} — \textquotedblleft{}shrine grounds; temple grounds\textquotedblright{}.

\begin{quote}
\textbf{\ja{けいだいで} \ja{こどもが} \ja{あそんでいます。}} — \emph{keidai de kodomo ga asonde imasu} — \textquotedblleft{}Children are playing in the shrine grounds.\textquotedblright{}
\end{quote}

\begin{cousinweb}
\textbf{\ja{けいだい}} (\emph{keidai}) is Sino-Japanese (kango): two Chinese characters read with their on'yomi, readings borrowed from Chinese. They are \emph{kei}, \textquotedblleft{}boundary\textquotedblright{} (Middle Chinese \emph{kjaengX}) and \emph{dai}, \textquotedblleft{}inside\textquotedblright{} (\emph{nwojH}) (KANJIDIC2; Middle Chinese in Baxter's transcription).
\end{cousinweb}

\begin{grammarlens}[title={What you've built}]
One more word for getting around.
\end{grammarlens}

\subsection*{Your turn}
\begin{itemize}
\raggedright
  \item \emph{Say it:} \emph{keidai}
  \item \emph{Say it:} \emph{keidai}, once more
  \item \emph{Say it:} \emph{keidai de kodomo ga asonde imasu} — Children are playing in the shrine grounds.
  \item \emph{Recall:} say the Japanese for a shrine gate, then the Japanese for to pass under, then say \emph{keidai} again
\end{itemize}

\begin{tcolorbox}[breakable,colback=teal!4,colframe=teal!35!black,title={Before you move on}]
What does \ja{けいだい} mean? (\textquotedblleft{}Shrine grounds; temple grounds\textquotedblright{}.) Say it once more.
\end{tcolorbox}
\section[\texorpdfstring{\ja{かんぬし} (kannushi)}{かんぬし (kannushi)}]{\ja{かんぬし} (kannushi) — a Shinto priest}

\label{lesson:JA-C406-kannushi}



\begin{quote}
Before the new one: say the Japanese for to pass under, then the Japanese for shrine grounds.
\end{quote}

\subsection*{You'll want to know: \ja{かんぬし}}
\textbf{\ja{かんぬし}} — \emph{kannushi} — \textquotedblleft{}a Shinto priest\textquotedblright{}.

\begin{quote}
\textbf{\ja{かんぬしが} \ja{まつりの} \ja{したくを} \ja{しています。}} — \emph{kannushi ga matsuri no shitaku o shite imasu} — \textquotedblleft{}The priest is getting ready for the festival.\textquotedblright{}
\end{quote}

\begin{cousinweb}
\textbf{\ja{かんぬし}} (\emph{kannushi}) is a native Japanese word (wago), inherited rather than borrowed from Chinese, built from \emph{kan}, \textquotedblleft{}god\textquotedblright{} (from \emph{kami}) and \emph{nushi}, \textquotedblleft{}master\textquotedblright{}; its kanji were attached later as kun'yomi (JMdict).
\end{cousinweb}

\begin{grammarlens}[title={What you've built}]
One more word for getting around.
\end{grammarlens}

\subsection*{Your turn}
\begin{itemize}
\raggedright
  \item \emph{Say it:} \emph{kannushi}
  \item \emph{Say it:} \emph{kannushi}, once more
  \item \emph{Say it:} \emph{kannushi ga matsuri no shitaku o shite imasu} — The priest is getting ready for the festival.
  \item \emph{Recall:} say the Japanese for to pass under, then the Japanese for shrine grounds, then say \emph{kannushi} again
\end{itemize}

\begin{tcolorbox}[breakable,colback=teal!4,colframe=teal!35!black,title={Before you move on}]
What does \ja{かんぬし} mean? (\textquotedblleft{}A Shinto priest\textquotedblright{}.) Say it once more.
\end{tcolorbox}
\section[\texorpdfstring{\ja{ささげる} (sasageru)}{ささげる (sasageru)}]{\ja{ささげる} (sasageru) — to offer up; to hold up}

\label{lesson:JA-C406-sasageru}



\begin{quote}
Before the new one: say the Japanese for shrine grounds, then the Japanese for a Shinto priest.
\end{quote}

\subsection*{You'll want to know: \ja{ささげる}}
\textbf{\ja{ささげる}} — \emph{sasageru} — \textquotedblleft{}to offer up; to hold up\textquotedblright{}.

\begin{quote}
\textbf{\ja{じんじゃに} \ja{はなを} \ja{ささげました。}} — \emph{jinja ni hana o sasagemashita} — \textquotedblleft{}We offered flowers at the shrine.\textquotedblright{}
\end{quote}

Polite \textbf{\ja{ささげます}}, \emph{sasagemasu}; an ichidan verb. Past (polite): \textbf{\ja{かんぬしが} \ja{はなを} \ja{ささげました。}}, \emph{kannushi ga hana o sasagemashita}.

\begin{cousinweb}
\textbf{\ja{ささげる}} (\emph{sasageru}) is a native Japanese word (wago), inherited rather than borrowed from Chinese; its kanji was attached later as a kun'yomi (JMdict).
\end{cousinweb}

\begin{grammarlens}[title={What you've built}]
One more word for getting around.
\end{grammarlens}

\subsection*{Your turn}
\begin{itemize}
\raggedright
  \item \emph{Say it:} \emph{sasageru}
  \item \emph{Say it:} \emph{sasageru}, once more
  \item \emph{Say it:} \emph{jinja ni hana o sasagemashita} — We offered flowers at the shrine.
  \item \emph{Recall:} say the Japanese for shrine grounds, then the Japanese for a Shinto priest, then say \emph{sasageru} again
\end{itemize}

\begin{tcolorbox}[breakable,colback=teal!4,colframe=teal!35!black,title={Before you move on}]
What does \ja{ささげる} mean? (\textquotedblleft{}To offer up; to hold up\textquotedblright{}.) Say it once more.
\end{tcolorbox}
